%%%%%%%%%%%%%%%%%%%%%%%%%%%%%%%%%%%%%%%%%%%%%%%%%%%%%%%%%%%%%%%%%%%%%%%%%%%
% Función: ejemplo de presupuesto de desarrollo de un prototipo.
%
% Uso: adapta los conceptos, cantidades y precios a tu trabajo.
%%%%%%%%%%%%%%%%%%%%%%%%%%%%%%%%%%%%%%%%%%%%%%%%%%%%%%%%%%%%%%%%%%%%%%%%%%%

\chapter{Presupuesto}
\label{cha:presupuesto}

Este ejemplo estima el coste de desarrollar un prototipo de robot móvil con
control por joystick. Incluye el diseño del hardware y del software, el
montaje y las pruebas. No estima el coste de fabricar una serie de robots ni
de instalar el sistema en su destino final.

Los precios, las tarifas y los porcentajes son supuestos del ejemplo, no
cotizaciones ni tipos que deban aplicarse a otros proyectos. Los importes
se expresan sin IVA hasta el resumen final.

\section{Costes directos}
\label{sec:presupuesto-directos}

\begin{samepage}
  \subsection{Costes de personal}
  \label{sec:presupuesto-mano}

  \begin{table}[H]
    \caption{Coste de personal.}
    \label{tab:recursos_humanos}
    \begin{tabularx}{\linewidth}{@{}>{\raggedright\arraybackslash}X r r r@{}}
      \hline
      \textbf{Perfil}                           & \textbf{Horas} & \textbf{Tarifa (€/h)} & \textbf{Importe (€)} \\
      \hline
      Ingeniero de Desarrollo Hardware (Junior) &            180 &                 30,00 &             5.400,00 \\
      Ingeniero de Desarrollo Software (Junior) &            120 &                 25,00 &             3.000,00 \\
      Ingeniero Senior / Director de Proyecto   &             35 &                 55,00 &             1.925,00 \\
      \hline
      \textbf{Total de personal}                &   \textbf{335} &                       &   \textbf{10.325,00} \\
      \hline
    \end{tabularx}
  \end{table}

  Las tarifas son referencias de facturación profesional sin IVA y se
  usan para valorar la dedicación en este ejemplo.

  Las 300 horas del estudiante (12 ECTS a razón de 25 horas por crédito)
  se reparten en 180 horas de desarrollo de hardware (diseño de esquemas y
  placas de circuito impreso, montaje y ensayos de laboratorio) y 120 horas
  de desarrollo de software (arquitectura del firmware y del código,
  programación, depuración y documentación). Las 35 horas de supervisión
  técnica y gestión del proyecto corresponden al ingeniero sénior y se
  añaden a las 300 horas del estudiante.
\end{samepage}
\pagebreak[0]

\begin{samepage}
  \subsection{Hardware y amortización}
  \label{sec:presupuesto-hardware}

  Los componentes que se incorporan al prototipo se imputan por su precio
  de compra. Las piezas impresas corresponden al juego montado en el robot.

  \nopagebreak
  \begin{table}[H]
    \caption{Componentes incorporados al prototipo.}
    \label{tab:recursos_hardware}
    \begin{tabularx}{\linewidth}{@{}>{\raggedright\arraybackslash}X r r r@{}}
      \hline
      \textbf{Componente}           & \textbf{Cantidad} & \textbf{Precio unitario (€)} & \textbf{Importe (€)} \\
      \hline
      Chasis                        &                 1 &                        50,00 &                50,00 \\
      Ruedas                        &                 4 &                        10,00 &                40,00 \\
      Baterías                      &                 2 &                        40,00 &                80,00 \\
      Motores con encoder           &                 2 &                        20,00 &                40,00 \\
      Piezas impresas en 3D (juego) &                 1 &                        15,00 &                15,00 \\
      Controladores de motor        &                 2 &                        30,00 &                60,00 \\
      Tarjetas LPC2129              &                 2 &                        50,00 &               100,00 \\
      Joystick analógico inductivo  &                 1 &                       110,00 &               110,00 \\
      \hline
      \textbf{Subtotal}             &                   &                              &      \textbf{495,00} \\
      \hline
    \end{tabularx}
  \end{table}
\end{samepage}
\pagebreak[0]

\begin{samepage}
  \medskip
  Para los equipos reutilizables se aplica una amortización lineal: se
  divide el precio de compra por la vida útil estimada en meses y se
  multiplica por los meses de uso imputados al proyecto.

  \nopagebreak
  \begin{table}[H]
    \caption{Amortización de los equipos reutilizables.}
    \label{tab:presupuesto-amortizacion}
    \begin{tabularx}{\linewidth}{@{}>{\raggedright\arraybackslash}X r r r r@{}}
      \hline
      \textbf{Equipo}       & \textbf{Compra (€)} & \textbf{Vida útil (meses)} & \textbf{Uso (meses)} & \textbf{Importe (€)} \\
      \hline
      Ordenador portátil    &            1.200,00 &                         48 &                    4 &               100,00 \\
      Equipo de laboratorio &            1.800,00 &                         60 &                    2 &                60,00 \\
      \hline
      \textbf{Subtotal}     &                     &                            &                      &      \textbf{160,00} \\
      \hline
    \end{tabularx}
  \end{table}
  \medskip
  El coste de hardware y amortización asciende a \textbf{655,00\,€}.
\end{samepage}
\pagebreak[0]

\begin{samepage}
  \subsection{Licencias de software}
  \label{sec:presupuesto-software}

  Se imputa la parte de cada licencia anual que corresponde a los meses
  de uso del proyecto. Este reparto es un criterio de imputación de costes,
  no implica que las licencias puedan contratarse por meses.

  \nopagebreak
  \begin{table}[H]
    \caption{Coste imputado de las licencias comerciales.}
    \label{tab:recursos_software}
    \begin{tabularx}{\linewidth}{@{}>{\raggedright\arraybackslash}X r r r@{}}
      \hline
      \textbf{Herramienta y uso}                         & \textbf{Precio (€/año)} & \textbf{Uso (meses)} & \textbf{Importe (€)} \\
      \hline
      Cadence OrCAD: diseño de circuitos                 &                  600,00 &                    4 &               200,00 \\
      MATLAB y Simulink: diseño y simulación del control &                  360,00 &                    2 &                60,00 \\
      Arm Keil MDK (µVision): programación y depuración  &                  240,00 &                    3 &                60,00 \\
      \hline
      \textbf{Subtotal}                                  &                         &                      &      \textbf{320,00} \\
      \hline
    \end{tabularx}
  \end{table}
\end{samepage}
\pagebreak[0]

\begin{samepage}
  \subsection{Material fungible}
  \label{sec:presupuesto-material}

  Se incluyen los materiales consumidos en el montaje y las pruebas, sin
  repetir las piezas que se han imputado al prototipo.

  \nopagebreak
  \begin{table}[H]
    \caption{Coste del material fungible.}
    \label{tab:presupuesto-fungible}
    \begin{tabularx}{\linewidth}{@{}>{\raggedright\arraybackslash}X r r r@{}}
      \hline
      \textbf{Concepto}                              & \textbf{Lotes} & \textbf{Precio por lote (€)} & \textbf{Importe (€)} \\
      \hline
      Resistencias, condensadores y LED para pruebas &              1 &                        60,00 &                60,00 \\
      Estaño y consumibles de montaje                &              1 &                        30,00 &                30,00 \\
      Material para montajes de ensayo               &              1 &                        30,00 &                30,00 \\
      \hline
      \textbf{Subtotal}                              &                &                              &      \textbf{120,00} \\
      \hline
    \end{tabularx}
  \end{table}
\end{samepage}
\pagebreak[0]

\begin{samepage}
  \section{Costes adicionales}

  \subsection{Costes indirectos y honorarios facultativos}
  \label{sec:presupuesto-costes-adicionales}

  Para estimar los gastos compartidos que no se han imputado de forma
  directa, como los suministros y el mantenimiento del espacio de trabajo,
  se aplica un 25\,\% al total de los costes directos.

  El ejemplo contempla una revisión técnica externa de la documentación y
  de los resultados de ensayo. Este servicio es independiente del trabajo
  del estudiante y de las 35 horas de supervisión y gestión ya imputadas
  a personal. Para estimarlo se aplica un 5\,\% a los costes directos.

  \nopagebreak
  \begin{table}[H]
    \caption{Cálculo de los costes indirectos y los honorarios facultativos.}
    \label{tab:presupuesto-costes-adicionales}
    \begin{tabularx}{\linewidth}{@{}>{\raggedright\arraybackslash}X r r r@{}}
      \hline
      \textbf{Concepto}         & \textbf{Base (€)} & \textbf{Porcentaje} & \textbf{Importe (€)} \\
      \hline
      Costes indirectos        &         11.420,00 &              25\,\% &             2.855,00 \\
      Honorarios facultativos  &         11.420,00 &               5\,\% &               571,00 \\
      \hline
    \end{tabularx}
  \end{table}
\end{samepage}
\pagebreak[0]

\begin{samepage}
  \section{Resumen del presupuesto}
  \label{sec:presupuesto-total}

  El resumen distingue los costes directos, los costes indirectos y los
  honorarios. A efectos de este ejemplo se aplica un IVA del 21\,\% al
  total antes de impuestos y se redondea la cuota al céntimo.

  \nopagebreak
  \begin{table}[H]
    \caption{Resumen económico del desarrollo del prototipo.}
    \label{tab:presupuesto_total}
    \begin{tabularx}{\linewidth}{@{}>{\raggedright\arraybackslash}X r@{}}
      \hline
      \textbf{Concepto}               & \textbf{Importe (€)} \\
      \hline
      Costes de personal              &            10.325,00 \\
      Hardware y amortización         &               655,00 \\
      Licencias de software           &               320,00 \\
      Material fungible               &               120,00 \\
      \hline
      \textbf{Costes directos}        &   \textbf{11.420,00} \\
      Costes indirectos (25\,\%)      &             2.855,00 \\
      Honorarios facultativos (5\,\%) &               571,00 \\
      \hline
      \textbf{Total sin IVA}          &   \textbf{14.846,00} \\
      IVA (21\,\%)                    &             3.117,66 \\
      \hline
      \textbf{Total con IVA}          &   \textbf{17.963,66} \\
      \hline
    \end{tabularx}
  \end{table}
\end{samepage}

%%% Local Variables:
%%% TeX-master: "../book"
%%% End:
